\chapter{The single-period binomial market}
\label{ch:binomial-one}

\begin{question}
In a market with a stock that goes up or down and a bond, when is there no free lunch, and what is the fair price of any claim?
\end{question}

\begin{whatfor}
\begin{itemize}
  \item Sets up a general discrete market: securities, trading strategies, portfolio value, self-financing, the no-arbitrage principle.
  \item In the binomial market, finds the cheapest hedge of a claim and shows its cost is an expected value under the risk-neutral probability $p^*=\frac{1+r-d}{u-d}$.
  \item Proves put--call parity.
\end{itemize}
\skipline{fs:binomial-one}
\end{whatfor}

\begin{prereq}
From \crefrange{ch:arbitrage}{ch:derivatives}: arbitrage, risk-neutral probability, call, put, contingent claim, bond, stock-and-bond portfolio. From Part~1: filtration, martingale, the binomial tree.
\end{prereq}

%=====================================================================
\section{A general discrete market}

\begin{setup}[Securities]
There are $d+1$ securities with prices $S_n=(S_n^0,S_n^1,\dots,S_n^d)$ at times $n=0,1,\dots,T$. Security $0$ is the bond, $S_n^0=(1+r)^n$; the others are risky.
\end{setup}

\begin{definition}[Trading strategy; value]
A \dfn{trading strategy} is a sequence $\Phi=(\Phi_n)$ of holdings $\Phi_n=(\Phi_n^0,\dots,\Phi_n^d)$, where $\Phi_n^i$ is the number of units of security $i$ held during period $n$ (from time $n-1$ to time $n$). It is chosen at time $n-1$, so $\Phi_n$ is $\F_{n-1}$-measurable. The \dfn{portfolio value} at time $n$ is
\[
  V_n(\Phi)=\Phi_n\cdot S_n=\sum_{i=0}^d\Phi_n^iS_n^i .
\]
\end{definition}

\begin{definition}[Self-financing]
A strategy is \dfn{self-financing} if rebalancing at time $n$ neither adds nor removes money:
\[
  \Phi_n\cdot S_n=\Phi_{n+1}\cdot S_n\qquad\text{for every }n .
\]
\end{definition}

\begin{result}[Gains of a self-financing strategy]
\label{res:gains}
For a self-financing strategy, the change of value over period $n+1$ is the holding times the price change:
\[
  V_{n+1}(\Phi)-V_n(\Phi)=\Phi_{n+1}\cdot(S_{n+1}-S_n),\qquad\text{so}\qquad V_n=V_0+\sum_{k=0}^{n-1}\Phi_{k+1}\cdot(S_{k+1}-S_k).
\]
\end{result}
\begin{proof}
$V_{n+1}-V_n=\Phi_{n+1}\cdot S_{n+1}-\Phi_n\cdot S_n=\Phi_{n+1}\cdot S_{n+1}-\Phi_{n+1}\cdot S_n$, using self-financing for the second term. Add up the periods.
\end{proof}

\begin{example}[A self-financing trade]
$S_0=(1,100)$, $S_1=(1.05,130)$ after an up-move. Holding $\Phi_1=(-30.48,0.4)$ (borrow $30.48$, own $0.4$ shares) is worth $\Phi_1\cdot S_1=-32+52=20$ at time $1$. Rebalancing to $\Phi_2=(-18.10,0.3)$ costs $\Phi_2\cdot S_1=-19.00+39=20$: the same, so the trade is self-financing (sell $0.1$ share for $13$, repay $12.38$ of the loan). If next $S_2=(1.1025,150)$, the gain is $\Phi_2\cdot(S_2-S_1)=-18.10(0.0525)+0.3(20)=5.05$.
\end{example}

\begin{definition}[Arbitrage; no-arbitrage principle]
An \dfn{arbitrage strategy} is a self-financing $\Phi$ with $V_0(\Phi)\le0$, $V_T(\Phi)\ge0$ and $P(V_T(\Phi)>0)>0$: it costs nothing, can never lose, and may win. The \dfn{no-arbitrage principle} says no such strategy exists.
\end{definition}

%=====================================================================
\section{The binomial market}

\begin{setup}[Binomial market]
A bond $B_t=(1+r)^t$ and a stock with $S_t=uS_{t-1}$ or $S_t=dS_{t-1}$, $t=1,\dots,T$, where $0<d<u$. A claim with expiry $T$ pays $X=g(S_T)$. This chapter takes $T=1$; Chapter~21 takes any $T$.
\end{setup}

\begin{result}[No arbitrage in the binomial market]
\label{res:binomial-na}
The binomial market is arbitrage-free if and only if $d<1+r<u$.
\end{result}
\begin{proof}
One direction here; Chapter~20 gives the other.
\begin{steps}
  \item If $1+r\le d$: borrow $S_0$ and buy one share. Today: $V_0=0$. At time $1$ the share is worth at least $dS_0\ge(1+r)S_0$, the debt: $V_1\ge0$, and $V_1>0$ after an up-move. Arbitrage.
  \item If $u\le1+r$: short one share and lend $S_0$. At time $1$ you have $(1+r)S_0\ge uS_0\ge S_1$: arbitrage again.
  \item So no arbitrage forces $d<1+r<u$. For the converse, Chapter~20 shows that a probability making the discounted stock a martingale rules out arbitrage, and $p^*$ below is such a probability. \qedhere
\end{steps}
\end{proof}

%=====================================================================
\section{Hedging a claim in one period}

A portfolio $(\Delta,b)$ holds $\Delta$ shares and $b$ bonds ($B_0=1$). Its value is $V_0=\Delta S_0+b$ today and $V_1=\Delta S_1+b(1+r)$ tomorrow. Write $X_u=g(uS_0)$, $X_d=g(dS_0)$.

\begin{definition}[Hedge; perfect hedge]
$(\Delta,b)$ is a \dfn{hedge} of $X$ if $V_1\ge X$ in both states:
\[
  \Delta uS_0+b(1+r)\ge X_u,\qquad \Delta dS_0+b(1+r)\ge X_d .
\]
It is a \dfn{perfect hedge} if both hold with equality.
\end{definition}

\begin{pic}
\begin{tikzpicture}
\begin{axis}[width=8.5cm,height=5.4cm,xmin=0,xmax=1.2,ymin=-0.7,ymax=0.7,xlabel={$\Delta$},ylabel={$b$},samples=2,domain=0:1.2]
  \addplot[draw=none,fill=thmc!12] coordinates {(0,0.6)(0.6,-0.24)(1.2,-0.48)(1.2,0.7)(0,0.7)} --cycle;
  \addplot[defc,very thick] {-1.4*x+0.6}; \node[defc,font=\scriptsize,anchor=west] at (axis cs:0.05,0.6) {(1): $b\ge-1.4\Delta+0.6$};
  \addplot[trapc,very thick] {-0.4*x}; \node[trapc,font=\scriptsize,anchor=north] at (axis cs:1.0,-0.42) {(2): $b\ge-0.4\Delta$};
  \addplot[only marks,mark=*] coordinates {(0.6,-0.24)};
  \node[font=\scriptsize,anchor=south west] at (axis cs:0.62,-0.22) {perfect hedge $(0.6,-0.24)$};
  \node[thmc,font=\scriptsize] at (axis cs:0.9,0.35) {all hedges};
\end{axis}
\end{tikzpicture}
\picnote{Hedges of the call in the example below ($S_0=1$, $u=1.75$, $d=0.5$, $r=0.25$, $K=1$). Each state gives a half-plane; the hedges are their intersection. The cheapest one, $V_0=\Delta S_0+b$ smallest, is the corner where both hold with equality.}
\end{pic}

\begin{theorem}[The perfect hedge and its price]
\label{thm:perfect-hedge}
If $d<1+r<u$, the perfect hedge is
\[
  \Delta=\frac{X_u-X_d}{(u-d)S_0},\qquad b=\frac{uX_d-dX_u}{(1+r)(u-d)},
\]
and its cost is
\[
  X^*=V_0=\frac{1}{1+r}\big[p^*X_u+(1-p^*)X_d\big]=\E_{P^*}\Big(\frac{X}{1+r}\Big),\qquad p^*=\frac{1+r-d}{u-d}\in(0,1).
\]
$X^*$ is the fair (no-arbitrage) price of $X$.
\end{theorem}
\begin{proof}
Solve the two equations, then rearrange the cost.
\begin{steps}
  \item Subtract the equations: $\Delta(u-d)S_0=X_u-X_d$. Substitute back: $b(1+r)=X_u-\Delta uS_0=\frac{(u-d)X_u-u(X_u-X_d)}{u-d}=\frac{uX_d-dX_u}{u-d}$.
  \item The cost:
  \[
    \begin{aligned}
    V_0=\Delta S_0+b&=\frac{X_u-X_d}{u-d}+\frac{uX_d-dX_u}{(1+r)(u-d)}\\
    &=\frac{1}{1+r}\Big[\frac{1+r-d}{u-d}X_u+\frac{u-(1+r)}{u-d}X_d\Big].
    \end{aligned}
  \]
  \item $d<1+r<u$ gives $0<p^*<1$. Any other price than $V_0$ lets you sell the dearer of (claim, hedge) and buy the cheaper: a sure profit. \qedhere
\end{steps}
\end{proof}

\begin{result}[The discounted stock is a martingale under $P^*$]
\label{res:disc-mg}
$\E_{P^*}\big(\frac{S_1}{1+r}\big)=S_0$, and no arbitrage holds exactly when $S_0$ lies strictly between $\frac{dS_0}{1+r}$ and $\frac{uS_0}{1+r}$.
\end{result}
\begin{proof}
$p^*uS_0+(1-p^*)dS_0=\frac{(1+r-d)u+(u-1-r)d}{u-d}S_0=(1+r)S_0$. And $S_0$ is a weighted average, with weights $p^*$, $1-p^*$ in $(0,1)$, of the two discounted prices exactly when it lies strictly between them.
\end{proof}

\begin{example}[A call and a put]
\label{ex:call-put}
$S_0=1$, $u=1.75$, $d=0.5$, $r=0.25$, $K=1$.
\begin{steps}
  \item $p^*=\frac{1.25-0.5}{1.75-0.5}=0.6$.
  \item Call: $X_u=0.75$, $X_d=0$. $C=\frac{1}{1.25}(0.6\cdot0.75)=0.36$. Hedge: $\Delta=\frac{0.75}{1.25}=0.6$, $b=\frac{0-0.5\cdot0.75}{1.25\cdot1.25}=-0.24$. Check: $0.6\cdot1.75-0.24\cdot1.25=0.75$ and $0.6\cdot0.5-0.3=0$.
  \item Put: $X_u=0$, $X_d=0.5$. $P=\frac{1}{1.25}(0.4\cdot0.5)=0.16$.
\end{steps}
\end{example}

%=====================================================================
\section{Put--call parity}

\begin{theorem}[Put--call parity]
\label{thm:parity}
For a European call and put on the same stock with the same strike $K$ and expiry $T$, in an arbitrage-free market,
\[
  S_0+P-C=\frac{K}{(1+r)^T}.
\]
\end{theorem}
\begin{proof}
Hold one share, one put, and sell one call.
\begin{steps}
  \item At $T$ it is worth $S_T+(S_T-K)^--(S_T-K)^+=S_T-(S_T-K)=K$ in every state, because $x^+-x^-=x$.
  \item A sure $K$ at time $T$ is worth $K(1+r)^{-T}$ today (a bond); any other price is an arbitrage. Today the holding costs $S_0+P-C$. \qedhere
\end{steps}
\end{proof}

In \cref{ex:call-put}: $P=\frac{1}{1.25}-1+0.36=0.16$, as computed directly.

%=====================================================================
\formulasheet{fs:binomial-one}

The one-period price, hedge and $p^*$ of this chapter are row~1 of Card~20.1 (beneath Chapter~20), which also covers markets with more outcomes or stocks.

\begin{mastercard}[label=card:selffin]{Value and gains of a trading strategy}
Holdings $\Phi_n$ are chosen at time $n-1$; prices $S_n$ include the bond $(1+r)^n$. A self-financing strategy never adds or takes out cash, so its value $V_n=\Phi_n\cdot S_n$ changes only through price moves of what it holds:
\[
  \boxed{\;V_n=V_0+\sum_{k<n}\Phi_{k+1}\cdot(S_{k+1}-S_k).\;}
\]
An arbitrage is a self-financing strategy with $V_0\le0$, $V_T\ge0$ and $P(V_T>0)>0$.
\begin{plugrows}
  \plugrow{Value of a holding}{\pstep{Lead: $V_n=\Phi_n\cdot S_n$, units times prices, summed.}\pstep{From the holding: $\Phi$ = (bonds, shares); $S$ = (bond price, share price).}\pstep{Multiply and add.}}{$V_n=\Phi_n\cdot S_n$}{all securities, bond included}{$\Phi=(-30.48,\,0.4)$, $S=(1.05,\,130)$: $-30.48(1.05)+0.4(130)=-32+52=20$}
  \plugrow{Is a trade self-financing?}{\pstep{Box: the value changes only through price moves, so rebalancing at time $n$ costs nothing.}\pstep{That is $\Phi_n\cdot S_n=\Phi_{n+1}\cdot S_n$: old and new holdings, same time-$n$ prices.}\pstep{Compute both sides (row 1).}}{equal $\Rightarrow$ self-financing}{at the time-$n$ prices}{old value $20$; new $\Phi_2=(-18.10,0.3)$ at $S_1=(1.05,130)$: $-18.10(1.05)+0.3(130)=-19+39=20$: equal, yes}
  \plugrow{Gain over a period}{\pstep{Box: $V_{n+1}-V_n$ is the term $k=n$ of the sum.}\pstep{That term: $\Phi_{n+1}\cdot(S_{n+1}-S_n)$, units held times price changes.}}{$V_{n+1}-V_n=\Phi_{n+1}\cdot(S_{n+1}-S_n)$}{self-financing}{$\Phi_2=(-18.10,0.3)$; bond $1.1025-1.05=0.0525$, share $150-130=20$: $-18.10(0.0525)+0.3(20)=-0.95+6=5.05$}
  \plugrow{Is it an arbitrage?}{\pstep{Lead: arbitrage = self-financing, $V_0\le0$, $V_T\ge0$, $P(V_T>0)>0$.}\pstep{Compute $V_0$ and $V_T$ in every outcome with row 1.}\pstep{Check all three.}}{yes / no}{all three needed}{$100\to200/50$, $r=0$, $K=150$, option at $20$; borrow $40$, buy 1 share, sell 3 options: $V_0=-40+100-60=0$; up: $-40+200-3(50)=10$; down: $-40+50-0=10$: an arbitrage}
  \plugrow{No arbitrage in the binomial market}{\pstep{Row 4 for one stock and the bond.}\pstep{If $1+r\le d$: borrow and buy the stock; it beats the loan in both states.}\pstep{If $u\le1+r$: short the stock and lend; the loan beats the stock in both states.}\pstep{If $d<1+r<u$: any holding that gains in one state loses in the other.}}{arbitrage-free $\iff d<1+r<u$}{one stock, two states, one period}{$u=1.1$, $d=0.95$, $r=0.1$: $1+r=1.1\ge u$: short $S_0$ of stock, lend it: $1.1S_0\ge S_1$ in both states}
\end{plugrows}
\end{mastercard}

\begin{mastercard}[label=card:parity]{Put--call parity}
Same stock, strike $K$, expiry $T$; $D_T$ = today's value of $1$ paid at $T$ ($(1+r)^{-T}$ in discrete time, $e^{-rT}$ in continuous time):
\[
  \boxed{\;S_0+P-C=K\,D_T,\qquad\text{at time }t:\ S_t+P_t-C_t=K\,D_{T-t}.\;}
\]
\begin{plugrows}
  \plugrow{Put from call}{\pstep{Box: $S_0+P-C=KD_T$.}\pstep{Solve for $P$.}\pstep{From the problem: $D_T=(1+r)^{-T}$, read $S_0$, $K$, $C$.}}{$P=C-S_0+KD_T$}{European options, no dividends}{$S_0=1$, $K=1$, $r=0.25$, $T=1$, $C=0.36$: $D_1=\frac1{1.25}=0.8$; $P=0.36-1+1(0.8)=0.16$}
  \plugrow{Call from put}{\pstep{Box: $S_0+P-C=KD_T$.}\pstep{Solve for $C$.}}{$C=P+S_0-KD_T$}{same}{$P=0.16$: $C=0.16+1-0.8=0.36$}
  \plugrow{Continuous time}{\pstep{Box with $D_T=e^{-rT}$.}\pstep{Rearrange: $C-P=S_0-KD_T$.}}{$C-P=S_0-Ke^{-rT}$}{European options, rate $r$ compounded continuously}{$S_0=80$, $K=85$, $r=0.05$, $T=1$: $85e^{-0.05}=80.85$; $C-P=80-80.85=-0.85$}
  \plugrow{Prices violate parity}{\pstep{Box: share $+$ put $-$ call pays $K$ at $T$ in every state, so it must cost $KD_T$.}\pstep{Compare the quoted $S_0+P-C$ with $KD_T$.}\pstep{Sell the dearer side, buy the cheaper; at $T$ the two positions cancel.}}{sure profit equal to the gap}{European options; any gap at all is an arbitrage}{$S_0+P-C=0.85$, $KD_T=0.8$: sell share and put, buy call (receive $0.85$), lend $0.8$; at $T$ the loan returns $K$ and covers the position: gain $0.85-0.8=0.05$}
  \plugrow{Straddle price}{\pstep{Straddle = call $+$ put (\cref{card:payoff}, row 4); prices add.}\pstep{Row 1: $P=C-S_0+KD_T$.}\pstep{Substitute.}}{$C+P=2C-S_0+KD_T$}{same $K$, $T$}{$C=0.36$: $2(0.36)-1+0.8=0.52$}
  \plugrow{Bounds on option prices}{\pstep{Box: $C=S_0+P-KD_T$, and $P\ge0$, so $C\ge S_0-KD_T$; also $C\ge0$.}\pstep{A call pays less than a share: $C\le S_0$.}\pstep{The same for the put: $P=C-S_0+KD_T\ge KD_T-S_0$; a put pays at most $K$: $P\le KD_T$.}}{$\max(S_0-KD_T,0)\le C\le S_0$ and $\max(KD_T-S_0,0)\le P\le KD_T$}{European options, no dividends}{$S_0=80$, $K=85$, $r=0.05$, $T=1$: $KD_T=80.85$: $\max(80-80.85,0)=0\le C\le80$; $\max(80.85-80,0)=0.85\le P\le80.85$}
  \plugrow{American call on a stock without dividends}{\pstep{Row 6 at time $t$: $C_t\ge S_t-KD_{T-t}$.}\pstep{With $r\ge0$, $D_{T-t}\le1$, so $S_t-KD_{T-t}\ge S_t-K$.}\pstep{Exercising now gives $S_t-K$; keeping the call is worth at least as much.}}{never exercise early: the American call costs the same as the European}{$r\ge0$, no dividends; false for puts, which can be worth exercising early}{$S_t=130$, $K=100$, $r=0.05$, $T-t=0.5$: $100e^{-0.025}=97.53$; $C_t\ge130-97.53=32.47>30=130-100$}
\end{plugrows}
\end{mastercard}
